\documentclass[12pt]{article}
\usepackage[utf8]{inputenc}
\usepackage[spanish]{babel}
\usepackage[a4paper, margin=2cm]{geometry}
\usepackage{tikz}
\usetikzlibrary{babel}
\usepackage[most]{tcolorbox}
\usepackage{amsmath}

\renewcommand{\familydefault}{\sfdefault}

\definecolor{medblue}{RGB}{43, 108, 176}
\definecolor{medteal}{RGB}{49, 151, 149}
\definecolor{medlight}{RGB}{235, 248, 255}
\definecolor{meddark}{RGB}{26, 32, 44}
\definecolor{angleorange}{RGB}{234, 88, 12}

\begin{document}
\pagestyle{empty}

\begin{tcolorbox}[
    enhanced, 
    colback=medlight!30, 
    colframe=medblue, 
    title={\Large \textbf{Ángulo Plano: Grados a Radianes}}, 
    halign title=center, 
    fonttitle=\bfseries, 
    arc=4mm, 
    boxrule=1.5pt,
    shadow={2mm}{-2mm}{0mm}{black!30!white}
]

    \vspace{0.2cm}
    \begin{tcolorbox}[colback=white, colframe=medteal, arc=2mm, boxrule=1.2pt, title=\textbf{Caso Clínico / Problema}]
        Un paciente en rehabilitación realiza un movimiento de flexión con su brazo describiendo un ángulo de \textbf{$60^\circ$}. Exprese la medida de este ángulo en radianes.
    \end{tcolorbox}

    \vspace{0.4cm}
    
    \begin{center}
    \begin{tikzpicture}[scale=1.5]
        % Líneas del ángulo
        \draw[line width=3pt, meddark, cap=round] (0,0) -- (4,0) node[right, font=\bfseries] {Eje Inicial};
        \draw[line width=3pt, meddark, cap=round] (0,0) -- (60:4) node[above right, font=\bfseries] {Flexión};
        
        % Sector circular (Ángulo)
        \filldraw[angleorange, fill opacity=0.2, draw=angleorange, line width=2pt] (0,0) -- (2,0) arc (0:60:2) -- cycle;
        
        % Etiqueta del ángulo
        \node[font=\Large\bfseries, text=angleorange!80!black] at (30:2.8) {$60^\circ$};
        
        % Puntos
        \filldraw[medblue] (0,0) circle (4pt) node[below left, font=\bfseries, text=meddark] {Articulación};
    \end{tikzpicture}
    \end{center}

    \vspace{0.4cm}
    \begin{tcolorbox}[colback=medteal!10, colframe=medteal, arc=2mm, boxrule=1.2pt, title=\textbf{Desarrollo Matemático y Solución}]
        Para convertir un ángulo de grados sexagesimales a radianes, multiplicamos el valor por el factor de conversión universal \textbf{$\left(\frac{\pi \text{ rad}}{180^\circ}\right)$}:
        \begin{align*}
            \theta &= 60^\circ \times \left( \frac{\pi \text{ rad}}{180^\circ} \right) \\[1ex]
            \theta &= \frac{60\pi}{180} \text{ rad} = \frac{\pi}{3} \text{ rad}
        \end{align*}
        Si deseamos el valor en formato decimal, considerando que $\pi \approx 3.1416$:
        \[ \theta \approx 1.047 \text{ rad} \]
        \textbf{Resultado:} El ángulo de flexión equivale a \textbf{$\frac{\pi}{3}$ radianes} o \textbf{$1.047$ rad}.
    \end{tcolorbox}
    
\end{tcolorbox}

\end{document}